\Needspace{8\baselineskip}\originalchapter{Vlasov 响应与 Landau 阻尼}\label{ch:landau}
\originalsection{流体闭合留下的问题}
流体纵波 $\omega^2\simeq\omega_p^2+3k^2T/m$ 可以解释热压修正, 却没有碰撞less阻尼. 原因在于一个平均速度不能记录与波相速接近的粒子. 对一维静态均匀分布 $f_{0s}(v)$ (积分为 $n_s$), 线性 Vlasov 和 Poisson 联立给
\begin{equation}
 \epsilon(\omega,k)=1-\sum_s\frac{\omega_{ps}^2}{k^2n_s}
 \int_{\mathcal L}\frac{f_{0s}'(v)}{v-\omega/k}\dd v=0.
\end{equation}
符号可用大相速展开核对: $\int f_0'/(v-z)\simeq n/z^2$, 因而首项为 $1-\omega_p^2/\omega^2$. 路径 $\mathcal L$ 不是任意选择的避极点方法, 而由 t=0 初值问题的单侧 Laplace 变换决定.

\originalsection{为什么极点下绕}
先在 $\operatorname{Im}\omega>0$ 解变换方程, 此时实 v 轴无奇点. 向下延拓求阻尼根时, 必须连续保持路径在移动极点下方. 对实相速, Sokhotski 公式给积分的主值加 $\ii\pi f_0'(\omega/k)$ (先取 k>0). 若在下半平面重新沿未变形实轴积分, 会丢掉初值因果性, 得到另一支函数. 衰减来自许多速度的相混合, 总 Vlasov 能量仍守恒, 单个波模的能量转入精细速度结构.

弱阻尼根写 $\omega=\omega_r+\ii\gamma$, $\epsilon_r(\omega_r)=0$, 则
$\gamma=-\epsilon_i(\omega_r)/[\partial_\omega\epsilon_r(\omega_r)]$.
单峰 Maxwell 在正相速有负斜率, 能吸收波能, 给 $\gamma<0$; 反转斜率的束流可放大波. 这个符号结论还要检查波能量分母, 不能把所有负斜率分布对所有模式都称稳定.

\originalsection{Maxwell 与 Lorentz 分布}
取 Maxwell 速度尺度 $v_t=\sqrt{2T/m}$, $\zeta=\omega/(kv_t)$, 定义因果 $Z(\zeta)=\pi^{-1/2}\int_{\mathcal L}\e^{-u^2}(u-\zeta)^{-1}\dd u$. 电子介电函数为 $1+2\omega_p^2(1+\zeta Z)/(k^2v_t^2)$. 展开 $Z=-1/\zeta-1/(2\zeta^3)-3/(4\zeta^5)+\cdots+\ii\sqrt\pi\e^{-\zeta^2}$, 恢复 $\omega_r^2\simeq\omega_p^2+3k^2T/m$ 和指数小的阻尼.

作业采用 Lorentzian $F(v)=v_e/[\pi(v^2+v_e^2)]$, 归一为1. 其二阶矩发散, $v_e$ 是宽度而非有限温度的热速. 残数计算并按因果延拓给
\begin{equation}
 \epsilon=1-\frac{\omega_p^2}{(\omega+\ii|k|v_e)^2},\qquad
 \omega=\pm\omega_p-\ii|k|v_e.
\end{equation}
这是一种可解析的相混合模型; 不可用其发散的温度硬套 Maxwell 压强闭合.

\originalsection{离子声波的电子屏蔽}
若 $|\omega|\ll|k|v_e$, Lorentz 电子响应近似 $1+\omega_{pe}^2/(k^2v_e^2)$. 加入离子宽度 $v_i$ 得
$\omega=\pm\omega_{pi}/\sqrt{1+\omega_{pe}^2/(k^2v_e^2)}-\ii|k|v_i$.
长波实相速为 $v_e\sqrt{m_e/m_i}$. 对 Maxwell 实际弱阻尼离子声波通常要求 $T_e\gg T_i$, 因为相速需比离子热速大; 两种分布的阻尼律不可互换.
